\section{Scope tests for the Fourier rigidity argument}
\label{sec:s5-scope-tests}

The central equality in MA-FQ-004 uses the actual completed divisor,
exact cancellation of the archimedean kernel, and the sign and
support of the remaining prime terms.  The following tests make
these hypotheses explicit.  They are not additional hypotheses in
that theorem and are not claims of first priority for the known
$L$-function or crystalline examples.

\paragraph{MA-FQ-005 (standard examples do not contradict the combined rigidity premises).}
The following statements hold.
\begin{enumerate}
\item The genuine degree-zero datum for $\zeta(s)/L(s,\chi_5)$,
where $\chi_5$ is the primitive even quadratic character modulo
five, has $K=0$ and nonnegative coefficients $f(n)$ for $n>1$.
Its full completed divisor is not supported on the critical line,
because it has poles at $s=0,1$.
\item The genuine degree-zero datum for
$L(s,\chi_5)/L(s,\chi_{13})$, with the two primitive even
quadratic characters, has infinitely many negative coefficients
at good primes.  Removing finitely many bad-prime progressions
does not give the sign hypothesis of MA-FQ-001.
\item The explicit integer-mass Fourier quasicrystal below is a
positive measure, but its atomic coefficient function is not
positive definite on $\mathbb R_d$.  Its ordinary Fourier
transform cannot be a nonnegative tempered measure plus a
finite-total-variation measure.  Its spectrum also contains a
frequency which is not the logarithm of a prime power.
\item The elementary one-prime function
$(1-p^{1/2-s})^{-1}$ has a central vertical pole lattice but
violates Booker's axiom (A1).  It is not an allowed bad-prime
correction in MA-FQ-001.
\end{enumerate}

\begin{proof}
\emph{The positive Euler ratio.}
Booker's actual Artin data for the trivial character and
$\chi_5$ have the same real gamma factor, so their difference
has $K=0$.  Logarithmic differentiation gives
\[
 f(n)=(1-\chi_5(n))\frac{\Lambda(n)}{\sqrt n}\ge0
                         \qquad(n>1).
 \tag{MA-FQ.30}
\]
The completed zeta function, before multiplying by any
pole-removing polynomial, has simple poles at $0$ and $1$.
The completed nontrivial Dirichlet function is nonzero at
both points: use its usual nonvanishing at $1$ and its
functional equation.  The completed quotient therefore
has these same two poles.  In height coordinates
$s=1/2+iz$, they are at $z=\pm i/2$, so the required real
support already fails, independently of the location of
any nontrivial zero.

Deleting those poles would change the explicit formula.
Indeed, for $h(z)=\int e^{ixz}g(x)\,dx$ their contribution is
\[
 h(i/2)+h(-i/2)=
                    \int_{\mathbb R}2\cosh(x/2)g(x)\,dx.
 \tag{MA-FQ.31}
\]
If the remaining signed divisor were real supported,
the formula for the Fourier transform of its negative
would contain the additional density $-2\cosh(x/2)$.
This is not the identity used in MA-FQ-001.
Its positive prime measure is not tempered either:
polynomial mass growth would imply convergence of every
positive Laplace transform, hence local convergence of
the Euler logarithm throughout $\operatorname{Re}s>1/2$.
Exponentiation would then remove the actual pole at $1$,
a contradiction.  Thus one cannot discard the endpoint
poles and retain the tempered positive-prime identity.

\emph{The ratio of two entire Dirichlet completions.}
Both characters are even, so the gamma kernels agree, and
their nontrivial completions are entire.  At an unramified
prime the difference datum has coefficient
\[
 f(p)=(\chi_5(p)-\chi_{13}(p))\frac{\log p}{\sqrt p}.
 \tag{MA-FQ.32}
\]
If $p\equiv3\pmod{65}$, the two character values are
$-1,+1$, respectively.  If $p\equiv11\pmod{65}$, they
are $+1,-1$.  These values follow directly from the
quadratic residues modulo five and thirteen.  Each of
the two reduced residue classes contains infinitely many
primes by the classical Dirichlet theorem.  In particular
no finite set of bad primes removes all the negative
coefficients.  This sign failure is independent of any
zero-location hypothesis.

\emph{The explicit crystalline test.}
Use Kurasov--Sarnak's polynomial
\[
 P(z_1,z_2)=1-\frac{z_1}{3}+\frac{z_2^2}{3}-z_1z_2^2,
 \qquad F(w)=P(2^{-w},3^{-w}).
 \tag{MA-FQ.33}
\]
Their Theorem~1 and equations (39), (41)--(43) construct
the positive crystalline zero measure of $F$.  Here is
a direct test of its relevant coefficient function.
Put
\[
 \xi_1=\log2,\quad \xi_2=\log3,\quad
 A=\xi_1/2+\xi_2,\quad D=\xi_1/2-\xi_2.
\]
Its real zero ordinates satisfy
\[
                 S(t)=3\sin(At)+\sin(Dt)=0.
 \tag{MA-FQ.34}
\]
One has $|D|<A$ and $A/D\notin\mathbb Q$, since a
rational ratio would make $\log2/\log3$ rational.
At a root, $|\sin(At)|\le1/3$, and hence
$|3A\cos(At)|\ge\sqrt8 A>|D|$.  Thus every root is
simple, including zero.  There is a nonzero root
$\lambda$ between $\pi/(2A)$ and $3\pi/(2A)$, because
$S$ is positive at the first endpoint and negative at
the second.

If $2\lambda$ were also a root, the double-angle formula
and $S(\lambda)=0$ would give
\[
           0=6\sin(A\lambda)
                     \{\cos(A\lambda)-\cos(D\lambda)\}.
 \tag{MA-FQ.35}
\]
If the sine vanishes, both $A\lambda$ and $D\lambda$
are integer multiples of $\pi$, contradicting
$A/D\notin\mathbb Q$.  Otherwise the cosines agree,
so $A\lambda\equiv D\lambda$ or $-D\lambda$ modulo
$2\pi$.  The corresponding equality or opposite
equality of the sines, together with $S(\lambda)=0$,
again makes both sines zero and gives the same
contradiction.  Therefore $2\lambda$ is not a root.

The atomic mass function $a$ of this measure consequently
has $a(0)=a(\lambda)=a(-\lambda)=1$ and
$a(2\lambda)=a(-2\lambda)=0$.  Its matrix on
$0,\lambda,2\lambda$ is not positive semidefinite:
\[
 (1,-1,1)
 \begin{pmatrix}1&1&0\\1&1&1\\0&1&1\end{pmatrix}
 (1,-1,1)^{\mathsf T}=-1.
 \tag{MA-FQ.36}
\]
Thus $a$ is not positive definite on $\mathbb R_d$.
If the ordinary Fourier transform of this atomic
measure were a nonnegative tempered measure plus a
finite measure, subtracting the latter's inverse
transform would be a bounded continuous correction.
The shrinking-bump proof in MA-FQ-001 would then make
its atomic coefficient function positive definite,
contradicting (MA-FQ.36).

There is also an independent Euler-support obstruction.
Set $v=z_2^2$.  Since
\[
 P(z_1,z_2)=(1+v/3)
       \left(1-z_1\frac{1/3+v}{1+v/3}\right),\qquad
 \frac{1/3+v}{1+v/3}
     =\frac13+\frac89v-\frac8{27}v^2+O(v^3),
\]
the coefficient of $z_1z_2^4$ in the actual Taylor
series for $\log P$ is $8/27$.
The Fourier summation formula therefore gives a
nonzero atom at $\xi_1+4\xi_2=\log162$.
The integer $162=2\cdot3^4$ is not a prime power.
This term cannot occur in the Euler logarithmic
derivative of the actual Artin/automorphic quotient
used in MA-FQ-001.

\emph{The central one-prime lattice.}
For the last example,
\[
 \log(1-p^{1/2-s})^{-1}
         =\sum_{j\ge1}\frac{p^{j(1/2-s)}}j.
\]
Comparison with Booker's normalization
$\sum_{n\ge2}(f(n)/\log n)n^{1/2-s}$
gives $f(p^j)=\log p$ for every $j$.
For any positive $N$, the sequence
$f(p^j)(j\log p)^N$ is unbounded.
This violates (A1).  Equivalently its reciprocal
frequency comb has nondecaying masses, rather than
the summable bad-prime tails of MA-FQ.1i.
This proves all four statements.
\end{proof}

\paragraph{Source locators and proof-check boundary.}
The Dirichlet examples use Booker's Definition~1.3,
author p.~3, and Examples~1.4(2)--(3), p.~4, in
\path{arXiv:1308.3067v2}; the first example is also
MA-LD-003.  Classical nonvanishing at one and
Dirichlet's theorem on primes in arithmetic progressions
are carried background inputs.  The exact character
values, endpoint correction, and A1 obstruction above
are checked directly.

Kurasov, P., \& Sarnak, P. (2020).
\emph{Stable polynomials and crystalline measures}.
Journal of Mathematical Physics, 61, 083501.
Author version \url{https://arxiv.org/abs/2004.05678}:
the summation formula and Theorem~1 with proof,
\S2, pp.~3--8; equations (39), (41)--(43),
\S\S2--3, pp.~8--10.  The complete invoked summation
argument and Theorem~1 proof were read in this run.
The matrix calculation above does not use the paper's
deeper additive-structure classification.

\paragraph{A checked family-hypothesis remark.}
Calegari's \S3 opening, publisher PDF p.~7, says that
three and five both split in the real quadratic field,
although the family hypotheses in Theorem~1.1 impose
only the condition at five.  This does not hide an
additional condition needed here: Sasaki's published
Theorem~1 requires complete splitting at five and
projective decomposition image of order two above
five, with no condition at three.  The stipulated
double-transposition Frobenius supplies those
requirements.  This is a statement/application check,
not a new proof of the modularity-lifting theorem:
Sasaki, S. (2013), \emph{On Artin representations and
nearly ordinary Hecke algebras over totally real fields},
Documenta Mathematica, 18, 997--1038,
Theorem~1, pp.~998--999,
\url{https://doi.org/10.4171/DM/419}.
